% ==========================================
% PHAN E - CHINH SACH BGP CHO AS A
% ==========================================

\section{Phần E -- Chính sách BGP cho AS A}

\subsection{Mục tiêu}

AS A có hai ISP: B nối với A2 qua L4
và C nối với A3 qua L5.
Nhóm cấu hình để lưu lượng đi ra ưu tiên ISP B,
dùng ISP C làm đường dự phòng;
đồng thời làm cho các AS khác ưu tiên
đi vào AS A qua ISP B.
AS A chỉ quảng bá LAN-A ra ngoài,
không làm transit giữa hai ISP.

\subsection{Triển khai cấu hình}

\subsubsection{Lưu lượng đi ra ưu tiên ISP B}

Nhóm dùng LOCAL\_PREF để ưu tiên
các route học từ B đi vào A2.
Các route này được đặt LOCAL\_PREF là 200,
cao hơn giá trị 100 của các route học từ C.
Thuộc tính này được trao đổi qua iBGP,
nên A3 cũng có thể chọn đường qua A2.

\textbf{Cấu hình trên A2:}

\begin{lstlisting}
conf t
route-map FROM-B permit 10
 set local-preference 200
exit

router bgp 64626
 neighbor 10.71.26.14 route-map FROM-B in
end
\end{lstlisting}

Ở đây, chính sách điều khiển lưu lượng
đi ra nhưng route-map lại được gắn theo hướng
\texttt{in}. Hai hướng này nói về hai việc khác nhau:
``đi ra'' là hướng gói tin từ AS A tới mạng ngoài,
còn \texttt{in} là hướng nhận thông tin BGP
từ B1 vào A2.
Nói đơn giản, để gửi gói tin ra qua B,
AS A ưu tiên các route nhận vào từ B.

Route-map \texttt{FROM-B} không có điều kiện
\texttt{match}, nên áp dụng LOCAL\_PREF 200
cho các route nhận từ neighbor này.
Khi đường qua B không còn hợp lệ,
BGP có thể chọn đường còn lại qua C.

\subsubsection{Lưu lượng đi vào ưu tiên qua ISP B}

Nhóm dùng AS-path prepend trên A3
để làm đường tới LAN-A trực tiếp qua C
bớt được ưu tiên.
Cụ thể, A3 thêm hai lần AS 64626
vào AS\_PATH khi quảng bá LAN-A sang C1:

\begin{lstlisting}
set as-path prepend 64626 64626
\end{lstlisting}

Ngoài AS cục bộ được thêm theo cơ chế eBGP,
hai lần prepend này làm AS\_PATH dài hơn.
Khi các thuộc tính ưu tiên trước AS\_PATH
không tạo ra khác biệt, C1 sẽ ưu tiên
đường tới LAN-A qua B thay vì trực tiếp qua A3.
Đây là cách tác động đến lựa chọn của AS bên ngoài;
kết quả trong mô hình được kiểm chứng
bằng traceroute ở phần sau.

\subsubsection{Chỉ quảng bá LAN-A, không làm transit}

Nhóm dùng prefix-list \texttt{ONLY-A}
để chỉ cho phép prefix LAN-A
\texttt{10.171.126.0/26}.
Lệnh \texttt{match ip address prefix-list}
gắn điều kiện này vào route-map xuất sang từng ISP.
Các prefix không khớp bị chặn bởi
implicit deny ở cuối route-map.

\textbf{Cấu hình trên A2, hướng tới ISP B:}

\begin{lstlisting}
conf t
ip prefix-list ONLY-A seq 5 permit 10.171.126.0/26

route-map TO-B permit 10
 match ip address prefix-list ONLY-A
exit

router bgp 64626
 neighbor 10.71.26.14 route-map TO-B out
end
\end{lstlisting}

\textbf{Cấu hình trên A3, hướng tới ISP C:}

\begin{lstlisting}
conf t
ip prefix-list ONLY-A seq 5 permit 10.171.126.0/26

route-map TO-C permit 10
 match ip address prefix-list ONLY-A
 set as-path prepend 64626 64626
exit

router bgp 64626
 neighbor 10.71.26.18 route-map TO-C out
end
\end{lstlisting}

Route-map \texttt{TO-C} ở trên là cấu hình cuối,
gồm cả lọc prefix và prepend.
Khi bổ sung điều kiện \texttt{match}
vào cùng sequence 10, lệnh \texttt{set}
đã có vẫn được giữ lại.
Nhờ vậy, A3 chỉ quảng bá LAN-A sang C
và vẫn làm dài AS\_PATH của route này.

AS A vẫn nhận các route từ hai ISP
để phục vụ lưu lượng của mình,
nhưng không quảng bá route của ISP này
sang ISP còn lại.
Do đó, hai ISP không học đường đi
qua AS A tới các mạng của nhau.

\subsection{Kiểm chứng kết quả}

\subsubsection{Bảng BGP trước và sau chính sách}

Nhóm dùng \texttt{show ip bgp} trên A2 và A3
để so sánh lựa chọn đường đi.

\partzeroimage
{images/configuration/part_E/A2_bgp_before.png}
{Bảng BGP trên A2 trước khi áp dụng chính sách}
{fig:partE-a2-before}

\partzeroimage
{images/configuration/part_E/A3_bgp_before.png}
{Bảng BGP trên A3 trước khi áp dụng chính sách}
{fig:partE-a3-before}

Trước khi áp dụng chính sách,
A2 chọn LAN-C \texttt{10.171.126.128/26}
qua A3, với AS\_PATH \texttt{64826},
thay vì đường qua B có AS\_PATH
\texttt{64726 64826}.
A3 cũng chọn đường trực tiếp tới C1.
Như vậy, đường ngắn hơn qua ISP C
đang được ưu tiên tới LAN-C.

\partzeroimage
{images/configuration/part_E/A2_bgp_after.png}
{Bảng BGP trên A2 sau khi áp dụng chính sách}
{fig:partE-a2-after}

\partzeroimage
{images/configuration/part_E/A3_bgp_after.png}
{Bảng BGP trên A3 sau khi áp dụng chính sách}
{fig:partE-a3-after}

Sau khi áp dụng chính sách, A2 chọn các route
tới LAN-B, LAN-C và LAN-D qua B1
\texttt{10.71.26.14}, với LOCAL\_PREF 200.
Trên A3, các route tương ứng qua A2
\texttt{10.171.26.2} cũng có LOCAL\_PREF 200
và được đánh dấu là đường tốt nhất.
Riêng LAN-C, đường qua C1 có LOCAL\_PREF 100
vẫn xuất hiện trên A3 như một lựa chọn dự phòng.
Điều này cho thấy LOCAL\_PREF cao hơn
đã khiến đường qua B được chọn
dù AS\_PATH dài hơn.

\subsubsection{Traceroute giữa LAN-A và LAN-C}

\begin{lstlisting}
! A1
traceroute 10.171.126.129 source 10.171.126.1

! C1
traceroute 10.171.126.1 source 10.171.126.129
\end{lstlisting}

\partzeroimage
{images/configuration/part_E/LAN_A_C_traceroute.png}
{Traceroute từ LAN-A tới LAN-C sau khi áp dụng chính sách}
{fig:partE-trace-a-c}

Các hop trong Hình \ref{fig:partE-trace-a-c}
lần lượt tương ứng với A2
(\texttt{10.71.26.2}), B1
(\texttt{10.71.26.14}), B3
(\texttt{10.71.26.30}) và C1
(\texttt{10.71.26.34}).
Lưu lượng từ LAN-A đi theo đường
A1 $\rightarrow$ A2 $\rightarrow$ B1
$\rightarrow$ B3 $\rightarrow$ C1,
xác nhận chiều đi ra sử dụng ISP B.

\partzeroimage
{images/configuration/part_E/LAN_C_A_traceroute.png}
{Traceroute từ LAN-C tới LAN-A sau khi áp dụng chính sách}
{fig:partE-trace-c-a}

Ở chiều ngược lại, các hop là B3
(\texttt{10.71.26.33}), B1
(\texttt{10.71.26.29}), A2
(\texttt{10.71.26.13}) và A1
(\texttt{10.171.126.1}).
Kết quả cho thấy C1 đi vào AS A qua ISP B,
phù hợp với mục tiêu của chính sách prepend.

\subsubsection{Kiểm tra các route quảng bá ra ISP}

\begin{lstlisting}
! A2
show ip bgp neighbors 10.71.26.14 advertised-routes

! A3
show ip bgp neighbors 10.71.26.18 advertised-routes
\end{lstlisting}

\partzeroimage
{images/configuration/part_E/A2_advertised_routes.png}
{Các prefix A2 quảng bá tới ISP B}
{fig:partE-a2-advertised}

\partzeroimage
{images/configuration/part_E/A3_advertised_routes.png}
{Các prefix A3 quảng bá tới ISP C}
{fig:partE-a3-advertised}

Cả hai output đều ghi \texttt{Total number of prefixes 1},
và prefix duy nhất là LAN-A
\texttt{10.171.126.0/26}.
Các mạng LAN-B, LAN-C và LAN-D
không được AS A quảng bá ra hai ISP,
phù hợp với yêu cầu không làm transit.

Hai ảnh này được dùng để kiểm chứng
danh sách prefix quảng bá.
Cột \texttt{Path} trong output hiện có
không thể hiện trực tiếp AS\_PATH sau prepend
mà C1 nhận được.
Vì vậy, phần kiểm chứng chiều đi vào
dựa trên cấu hình \texttt{TO-C}
và traceroute từ C1 tới LAN-A,
không suy ra AS\_PATH đã prepend từ hai ảnh trên.

\newpage
